\section{Conclusion}
\label{sec:conclusion}

The ready-cohort boundary turns GPU agent control into a falsifiable systems
question. A workload must supply enough same-group events before their launch
deadlines, and the control path must account for where intermediate decisions
are observed. Under the frozen trace model, exact sliding-deadline packing
raises eligible share from \TraceFixedPct\% to \TraceExactPct\%, recovering
\TraceGapClosurePct\% of the opportunity lost by fixed windows. The full
surface also shows where cohort supply collapses to zero.

The mechanism study establishes the second gate. Keeping the tested binary
decision on device reduces cohort-horizon wall time in all
\ResidentCells{} cells across four named placements, with ratios from
\ResidentRatioMin$\times$ to \ResidentRatioMax$\times$. A nested device graph
that removes no host decision loses in all \NativeCells{} cells across five
placements. Device launch alone is therefore insufficient; the gain appears
only in the treatment that removes the matched host-mediated decision epoch.

These results provide the workload budget and the mechanism test for an online
route compactor. Its decisive measurements are achieved share $A$ relative to
$F$ and $\pstar$, raw P99, CPU core-seconds, exact effects, task utility, cost,
and shared-inference guardrails. If it cannot convert offline headroom into
online work or displace CPU without harming model service, the boundary rejects
the GPU design. If it can, the same quantities show exactly where GPU-resident
agent control belongs in a datacenter.
